\chapter{Estado da Arte}

Antes de se justificar o desenvolvimento de um modelo físico e dinâmico de árvore brônquica, importa perceber onde se situa este trabalho relativamente ao que já existe em simulação médica para broncoscopia, pelo que este capítulo percorre o estado atual do conhecimento nesta área, desde o papel da simulação na aprendizagem da broncoscopia e o conceito de fidelidade que permite comparar diferentes simuladores entre si, até ao contributo da impressão 3D na criação de modelos anatómicos e aos principais movimentos fisiológicos observáveis durante o procedimento, que servem de referência direta ao trabalho desenvolvido nos capítulos seguintes.

\section{Broncoscopia e treino por simulação}
A broncoscopia flexível é um procedimento médico endoscópico que permite observar diretamente a árvore traqueobrônquica, recolher amostras, auxiliar diagnósticos e apoiar intervenções terapêuticas. A aprendizagem deste procedimento exige coordenação entre visão broncoscópica, manipulação fina do broncoscópio e interpretação anatómica. A simulação pormenorizada é, por isso, uma ferramenta relevante para repetir tarefas, treinar orientação espacial e reduzir a dependência inicial de treino em doentes reais \parencite{Davoudi2008,Kennedy2013,Fielding2014,Nilsson2017}.

As revisões e estudos de treino em broncoscopia apontam para ganhos na aquisição de competências quando os simuladores são integrados de forma estruturada no ensino. O valor pedagógico de um simulador depende não só do realismo visual, mas também da relação entre o exercício proposto, os objetivos de aprendizagem e a forma como o desempenho é avaliado \parencite{Kennedy2013,Fielding2014,Gopal2018}.

Trabalhos mais recentes procuram comparar diretamente diferentes modalidades de treino entre si. Um estudo dinamarquês comparou quatro modelos distintos de treino em broncoscopia, incluindo modelos físicos e virtuais, concluindo que os modelos físicos de baixo custo se mantêm competitivos face a alternativas tecnologicamente mais sofisticadas e dispendiosas \parencite{Kronborg2024}, enquanto outro estudo mostrou que sessões curtas de treino guiadas por um simulador produzem resultados de aprendizagem semelhantes aos obtidos sob orientação direta de um tutor humano, reforçando o valor pedagógico da simulação mesmo sem supervisão constante \parencite{Schertel2021}.

\begin{figure}[htbp]
  \centering
  \includegraphics[width=0.86\textwidth]{artigos/modelos_treino_broncoscopia.png}
  \caption{Modelos de treino em broncoscopia \parencite{Kronborg2024}.}
  \label{fig:modelos-treino-broncoscopia}
\end{figure}

\section{Fidelidade dos simuladores}
A fidelidade de um simulador pode ser entendida em várias dimensões: 
\begin{itemize}
    \item Fidelidade anatómica $\rightarrow$ refere-se à correspondência entre a geometria do modelo e a anatomia real, como à forma, proporção e organização espacial das vias aéreas;
    \item Fidelidade mecânica $\rightarrow$ refere-se à resistência ao avanço, flexibilidade e resposta ao contacto; 
    \item Fidelidade funcional $\rightarrow$ refere-se à capacidade de reproduzir movimentos, variações de calibre e eventos transitórios como a tosse.
\end{itemize}

\begin{figure}[htbp]
  \centering
  \begin{tikzpicture}[
    axis/.style={draw, rounded corners=4pt, align=center, minimum width=3.8cm, minimum height=1.2cm, font=\small, thick},
    arr/.style={-Latex, thick}
  ]
    \node[axis, fill=blue!8] (anat) {Fidelidade\\anatómica};
    \node[axis, fill=green!10, right=1.0cm of anat] (mec) {Fidelidade\\mecânica};
    \node[axis, fill=orange!12, right=1.0cm of mec] (func) {Fidelidade\\funcional};
    \node[draw, rounded corners=5pt, align=center, below=1.0cm of mec, minimum width=5.0cm, minimum height=1.0cm, fill=gray!10] (treino) {Valor pedagógico\\do simulador};
    \draw[arr] (anat.south) -- (treino.north west);
    \draw[arr] (mec.south) -- (treino.north);
    \draw[arr] (func.south) -- (treino.north east);
  \end{tikzpicture}
  \caption{Dimensões de fidelidade consideradas num simulador físico.}
  \label{fig:dimensoes-fidelidade}
\end{figure}

Esta distinção aproxima-se da lógica proposta por Miller para a avaliação de competências clínicas, que separa aquilo que um profissional sabe, sabe fazer, demonstra saber fazer e efetivamente faz na prática, situando os simuladores físicos precisamente na transição entre o demonstrar e o fazer \parencite{Witheridge2019}. Importa, no entanto, notar que a relação entre fidelidade e aprendizagem não é necessariamente linear, havendo trabalhos recentes que questionam se um maior realismo se traduz sempre em mais competência adquirida, sugerindo que o envolvimento ativo do formando pode pesar tanto quanto o rigor da réplica anatómica \parencite{Pico2026}.

Na broncoscopia, um modelo estático pode ser útil para situações de treinos iniciais de navegação e reconhecimento anatómico. No entanto, durante um procedimento real, o operador encontra-se num ambiente vivo e dinâmico, em constante movimento. A respiração, os batimentos cardíacos, a tosse, a sedação e a própria interação do broncoscópio com as paredes internas exercem comportamentos e reações que modificam continuamente o ambiente e a sensação de controlo. Esta diferença cria oportunidade para os simuladores físicos de baixo custo que incorporem movimentos e fisiologia de forma controlada \parencite{Davoudi2008,Nilsson2017,Fu2023}.

\section{Simuladores de broncoscopia}
Os atuais simuladores de broncoscopia recorrem a vários tipos de sistemas, nomeadamente treinos em ambientes puramente virtuais, em manequins comerciais, em modelos de bancada e, mais recentemente, em modelos impressos em 3D.

Os sistemas virtuais, como o simulador ORSIM, permitem gerar cenários de dificuldade variável e obter métricas automáticas de desempenho, tendo demonstrado boa validade e fiabilidade na avaliação de competências em broncoscopia flexível, distinguindo de forma consistente entre principiantes e peritos \parencite{Baker2016}. Contudo, estes sistemas dependem de hardware háptico dedicado para simular a resistência ao avanço e o contacto com as paredes das vias aéreas, um aspeto que continua a ser tecnicamente difícil de reproduzir com fidelidade, uma vez que a deteção de força na interação entre instrumento e tecido em endoscopia minimamente invasiva (termo mais abrangente, que inclui a broncoscopia como um dos seus procedimentos) permanece uma área ativa de investigação, sem uma solução consensual e amplamente adotada \parencite{Hosseinabadi2021}. Modelos físicos baseados em tecido biológico real, como um biosimulador reutilizável construído a partir de pulmões suínos ventilados, procuram contornar esta limitação ao oferecer uma resposta tátil genuína, tendo sido avaliados favoravelmente por broncoscopistas experientes e principiantes em comparação direta com um simulador virtual comercial \parencite{Garner2019}. Estes modelos biológicos, no entanto, não são reutilizáveis indefinidamente, exigem preservação especializada e apresentam variabilidade anatómica de espécime para espécime, o que limita a sua padronização enquanto ferramenta de ensino e de avaliação objetiva.

Esta lacuna, entre a normalização própria dos sistemas virtuais e o realismo tátil dos modelos biológicos, abre espaço para os modelos impressos em 3D, que procuram oferecer geometria anatómica consistente e reprodutível, associada a materiais capazes de aproximar a resposta mecânica dos tecidos reais.

\begin{figure}[H]
  \centering
  \includegraphics[width=0.74\textwidth]{artigos/broncoscopia_virtual_dunnbeltran.png}
  \caption{Exemplo de simulação virtual em broncoscopia \parencite{DunnBeltran2026}.}
  \label{fig:broncoscopia-virtual-dunnbeltran}
\end{figure}

\begin{figure}[H]
  \centering
  \includegraphics[width=0.76\textwidth]{artigos/simulador_biologico_garner.png}
  \caption{Biossimulador para treino de broncoscopia \parencite{Garner2019}.}
  \label{fig:simulador-biologico-garner}
\end{figure}

\section{Impressão 3D aplicada à simulação médica}
A impressão 3D é uma tecnologia de produção de objetos tridimensionais, realizada através da adição de material camada a camada, podendo o método de adição ou aglomeração dos materiais variar bastante consoante a tecnologia utilizada, entre as quais se destacam o FDM, a SLA e a SLS, cada uma assente num princípio de fabrico distinto e, por isso, mais adequada a diferentes tipos de aplicação.
Com o passar dos anos e o amadurecimento da tecnologia e dos materiais disponíveis, este método de produção veio proporcionar uma oportunidade inovadora na criação de estruturas e geometrias complexas, a um baixo custo e numa fração do tempo face aos outros métodos de produção, afetando todas as áreas científicas e industriais, incluindo a medicina. Esta permite a produção de protótipos anatómicos com custo relativamente baixo, a correção e o ajuste de dimensões em pouco tempo e o teste de materiais inovadores que, com outros métodos, seria impensável, tanto pelo custo de produção como pela complexidade dos processos.

\subsection{Técnicas de impressão 3D}
\label{subsec:tecnicas-impressao-3d}
O FDM (Fused Deposition Modeling, ou modelação por deposição de material fundido) é a tecnologia mais acessível e mais difundida, consistindo na extrusão de um filamento termoplástico aquecido através de um bico, que deposita o material camada a camada sobre uma base de impressão. É compatível com uma grande variedade de materiais, como o PLA, o PETG, o ABS ou filamentos flexíveis como o TPU, permitindo a produção de peças estruturalmente resistentes a um custo reduzido, ainda que com uma resolução e um acabamento de superfície mais limitados, sendo geralmente visíveis as linhas de cada camada. Para geometrias com zonas em consola ou sem apoio inferior, o FDM necessita ainda de estruturas de suporte, que têm de ser removidas após a impressão.

A SLA (Stereolithography) recorre a uma cuba de resina fotopolimérica líquida, curada seletivamente camada a camada por uma fonte de luz ultravioleta, um laser ou um projetor. Esta tecnologia permite obter peças com elevada resolução e um acabamento de superfície muito mais fino do que o FDM, sendo por isso indicada para modelos com geometrias detalhadas. Em contrapartida, exige equipamento de pós-processamento dedicado, nomeadamente para lavagem das peças em álcool isopropílico e cura adicional sob luz ultravioleta, além de recorrer a materiais mais dispendiosos e, em geral, mais frágeis do que os filamentos termoplásticos.

Por fim, a SLS (Selective Laser Sintering) utiliza um laser para sinterizar seletivamente um material em pó, tipicamente uma poliamida, camada a camada dentro de um leito de pó. Uma das suas principais vantagens é dispensar estruturas de suporte dedicadas, uma vez que o próprio pó não sinterizado sustenta as zonas em consola durante a impressão, permitindo fabricar geometrias mais complexas. As peças resultantes apresentam ainda propriedades mecânicas mais uniformes do que as obtidas por FDM. Contudo, o equipamento e os materiais são substancialmente mais caros, o que torna esta tecnologia mais comum em contexto industrial do que em pequenos projetos ou protótipos.

No contexto deste projeto, esta comparação foi particularmente relevante na escolha entre FDM e SLA para a impressão da árvore brônquica, decisão discutida em detalhe na \cref{sec:disposicao-impressao}.

A utilização clínica e educativa da impressão 3D em modelos anatómicos exige, contudo, um controlo de qualidade rigoroso ao longo de todo o processo, desde a segmentação da imagem médica até à impressão final, sendo possível identificar e quantificar erros próprios de cada etapa, como o erro de segmentação, o erro de edição digital e o erro de impressão propriamente dito, cuja soma determina a fidelidade geométrica final do modelo obtido \parencite{Schulze2024}. Esta preocupação com o rigor dimensional não se limita a modelos de treino, estendendo-se também a aplicações clínicas diretas, como o desenho e o fabrico de stents personalizados para o tratamento de estenoses das vias aéreas, benignas ou malignas, a partir da anatomia específica de cada doente \parencite{Aravena2023,Yilmaz2022,Shan2019}, ou ainda à produção de modelos de traqueia para treino de cricotiroidotomia de emergência e de modelos de vias aéreas superiores para treino de entubação endotraqueal, ambos criados com o objetivo de tornar este tipo de treino mais acessível e menos dependente de equipamento comercial dispendioso \parencite{Doucet2017,Park2023}.

\begin{figure}[htbp]
  \centering
  \includegraphics[width=0.92\textwidth]{artigos/workflow_modelos_3d.png}
  \caption{Fluxo de produção e validação de modelos anatómicos impressos em 3D \parencite{Schulze2024}.}
  \label{fig:workflow-modelos-3d}
\end{figure}

\subsection{Aplicação em simuladores de árvores brônquicas}
No caso dos simuladores de árvores brônquicas, os modelos impressos em 3D ganharam uma grande importância por permitirem transformar dados anatómicos, obtidos por exemplo a partir de exames de tomografia computorizada (TAC), em geometrias físicas. Em broncoscopia, esta abordagem pode representar a árvore traqueobrônquica com boa correspondência espacial e permitir a criação de simuladores específicos para treino de navegação \parencite{Pedersen2017,Fu2023}, ainda que a sua principal limitação, quando usados de forma isolada, seja a ausência de movimento e de resposta fisiológica. Modelos multi-material, combinando simultaneamente componentes rígidos e flexíveis numa só impressão, têm vindo a ser explorados precisamente para aproximar a experiência tátil da broncoscopia real sem comprometer a robustez estrutural do conjunto \parencite{Ho2019}, enquanto outras abordagens recorrem à impressão a cores, atribuindo uma cor distinta a cada segmento ou ramo brônquico, de modo a tornar mais percetível a organização da árvore traqueobrônquica e a facilitar simultaneamente o treino de navegação e a aprendizagem da nomenclatura anatómica \parencite{Leba2023}.

\begin{figure}[htbp]
  \centering
  \includegraphics[width=0.74\textwidth]{artigos/modelos_multimaterial_broncoscopia.png}
  \caption{Modelos brônquicos impressos em multi-material \parencite{Ho2019}.}
  \label{fig:modelos-multimaterial-broncoscopia}
\end{figure}

Além dos modelos multi-material, existem abordagens que exploram a cor e o planeamento personalizado para tornar a anatomia mais legível ou adaptar o modelo a uma necessidade clínica concreta.

\begin{figure}[H]
  \centering
  \includegraphics[width=0.68\textwidth]{artigos/modelo_vias_colorizado.png}
  \caption{Modelo de vias aéreas segmentado por cor \parencite{Leba2023}.}
  \label{fig:modelo-vias-colorizado}
\end{figure}

Contudo, a ideia de impressão destes modelos apresenta desafios e detalhes muito importantes, tais como a escolha do material, pois facilmente influencia o grau de dificuldade de impressão, a flexibilidade, a resistência, a sensação ao toque e a disponibilidade de cor. Materiais rígidos podem servir para validar forma e montagem, enquanto materiais flexíveis podem aproximar-se melhor à resposta mecânica, ainda que aumentem consideravelmente a dificuldade de impressão e calibração. Assim, a impressão 3D deve ser entendida também como uma plataforma de iteração, não apenas como uma técnica de fabrico final.

\begin{figure}[H]
  \centering
  \includegraphics[width=0.62\textwidth]{artigos/stent_aereo_personalizado.png}
  \caption{Planeamento de stent personalizado para vias aéreas \parencite{Aravena2023}.}
  \label{fig:stent-aereo-personalizado}
\end{figure}

\section{Movimentos observáveis durante a broncoscopia}
A árvore traqueobrônquica não se comporta como uma estrutura fixa durante uma broncoscopia. Mesmo quando o broncoscópio se mantém praticamente parado, a imagem observada pode apresentar deslocações lentas, alterações de orientação, variações do lúmen e perturbações rápidas. 
Para o operador, estes movimentos não aparecem isolados: 
\begin{itemize}
    \item A respiração cria uma componente periódica dominante;
    \item O batimento cardíaco acrescenta pequenas oscilações locais;
    \item A tosse introduz eventos abruptos;
    \item O contacto do broncoscópio com a parede pode modificar temporariamente a posição ou a forma do segmento observado.
\end{itemize}

Do ponto de vista de um simulador físico, estes fenómenos não precisam de ser reproduzidos com toda a complexidade fisiológica para terem valor pedagógico, mesmo que acabem sempre por apresentar um maior realismo. O mais importante é que o utilizador deixe de observar um modelo completamente estático e passe a treinar num ambiente com movimento por vezes imprevisível e suficientemente coerente com o que é visto durante o procedimento real. Por isso, os movimentos descritos de seguida foram analisados quanto à sua relevância visual, dificuldade de implementação e utilidade, seguindo a categorização e os intervalos de frequência e amplitude propostos num documento de trabalho elaborado especificamente para este projeto pela equipa de Pneumologia do Centro Hospitalar Universitário do Algarve \parencite{Santos2026}, que serviu de base clínica direta às secções seguintes.

\subsection{Movimento respiratório}
O movimento respiratório é a componente mais frequente e previsível, ocorrendo durante o processo cíclico da inspiração e expiração, descrito como a deslocação dos pulmões com predominância craniocaudal, acompanhado por pequenas componentes laterais e por alterações locais da geometria das vias aéreas. A frequência respiratória normal situa-se aproximadamente entre 12 e 20 ciclos por minuto, embora possa variar com ansiedade, sedação, esforço respiratório, patologia e contexto clínico.

Análises mostram que a fase respiratória pode alterar a posição de estruturas pulmonares e influenciar modelos de fluxo nas vias aéreas \parencite{Chen2015,Gunatilaka2020}, existindo ainda, além da deslocação global, variações do calibre brônquico ao longo do ciclo respiratório, um fenómeno já descrito por trabalhos clássicos que demonstraram os mecanismos das alterações de calibre dos brônquios durante a respiração normal e das suas variações rítmicas \parencite{Heinbecker1927,Ellis1936}. Mais recentemente, técnicas de modelação computacional têm procurado prever e simular este movimento de forma específica ao paciente, recorrendo a modelos de elementos finitos capazes de estimar o deslocamento dos brônquios ao longo do ciclo respiratório a partir de imagem médica, alcançando erros de previsão inferiores a 3\% quando validados com dados reais \parencite{Kim2023}. Para o protótipo desenvolvido, esta componente foi interpretada como um movimento lento e periódico, adequado para ser representado por trajetórias sinusoidais, suaves e repetíveis.

\begin{figure}[htbp]
  \centering
  \includegraphics[width=0.62\textwidth]{fotos/movimento_respiratorio_arvore.png}
  \caption{Representação esquemática do movimento respiratório da árvore brônquica.}
  \label{fig:movimento-respiratorio-arvore-estado-arte}
\end{figure}

\begin{figure}[htbp]
  \centering
  \includegraphics[width=0.88\textwidth]{artigos/modelo_movimento_respiratorio_kim.png}
  \caption{Modelo computacional do movimento respiratório \parencite{Kim2023}.}
  \label{fig:modelo-movimento-respiratorio-kim}
\end{figure}

\subsection{Movimento cardíaco}
O movimento cardíaco introduz oscilações mais rápidas e de menor amplitude, com intervalos regulares entre ocorrências, cuja influência tende a ser mais percetível em regiões próximas do coração, em particular junto ao brônquio principal esquerdo e a estruturas adjacentes. Ao contrário da respiração, que domina a posição global da árvore brônquica, o batimento cardíaco surge antes como uma vibração ou perturbação sobreposta.

Estudos com imagem 4D confirmam esta distinção, tendo quantificado separadamente as contribuições respiratória e cardíaca no movimento de estruturas pulmonares e observado uma influência cardíaca proporcionalmente maior em lesões centrais e superiores, mais próximas do coração e dos grandes vasos \parencite{Savanovic2023}. Esta influência não se limita à posição da árvore brônquica, estendendo-se também ao próprio fluxo de ar no seu interior. Imagem dinâmica obtida por radiação de sincrotrão revelou que o batimento cardíaco gera oscilações de pressão e de fluxo de ar mensuráveis nos brônquios principais, num fenómeno localizado junto à carina até então difícil de demonstrar experimentalmente \parencite{Dubsky2018}. Em casos anatómicos particulares, esta proximidade pode mesmo traduzir-se em compressão mecânica direta do brônquio principal esquerdo por estruturas cardiovasculares adjacentes, como a aorta descendente ou uma artéria pulmonar dilatada, um efeito documentado em população pediátrica com anomalias intracardíacas \parencite{Ando2024}.

\begin{figure}[htbp]
  \centering
  \begin{tikzpicture}[
    x=0.9cm,
    y=1.0cm,
    axis/.style={gray!55, thin},
    ecg/.style={red!70!black, very thick, line join=round, line cap=round},
    lab/.style={font=\small}
  ]
    \foreach \x in {0,...,12} {
      \draw[axis] (\x,-0.8) -- (\x,1.7);
    }
    \foreach \y in {-0.5,0,0.5,1,1.5} {
      \draw[axis] (0,\y) -- (12,\y);
    }
    \draw[ecg]
      (0,0) -- (0.6,0) .. controls (0.8,0.15) and (1.0,0.15) .. (1.2,0)
      -- (1.7,0) -- (1.9,-0.35) -- (2.1,1.55) -- (2.3,-0.55) -- (2.55,0)
      -- (3.2,0) .. controls (3.6,0.45) and (4.0,0.45) .. (4.45,0)
      -- (5.4,0)
      .. controls (5.6,0.12) and (5.85,0.12) .. (6.05,0)
      -- (6.55,0) -- (6.75,-0.35) -- (6.95,1.45) -- (7.15,-0.5) -- (7.42,0)
      -- (8.1,0) .. controls (8.5,0.42) and (8.9,0.42) .. (9.35,0)
      -- (10.1,0)
      .. controls (10.3,0.12) and (10.55,0.12) .. (10.75,0)
      -- (11.25,0) -- (11.45,-0.3) -- (11.65,1.35) -- (11.82,-0.45) -- (12,0);
    \node[lab, red!70!black] at (2.1,1.85) {QRS};
    \node[lab] at (1.0,0.42) {P};
    \node[lab] at (3.8,0.72) {T};
    \draw[-Latex, thick] (2.1,-0.95) -- (2.1,-0.62);
    \draw[-Latex, thick] (6.95,-0.95) -- (6.95,-0.58);
    \draw[Latex-Latex, thick] (2.1,-1.12) -- node[below, font=\small] {período entre batimentos} (6.95,-1.12);
  \end{tikzpicture}
  \caption{Representação esquemática de um ciclo cardíaco através da onda de ECG.}
  \label{fig:ecg-batimento-cardiaco}
\end{figure}

Num modelo de treino, esta componente é útil para quebrar a sensação de movimento puramente sinusoidal e para acrescentar uma perturbação de menor amplitude, pelo que, no contexto deste projeto, o batimento cardíaco não foi tratado como uma reconstrução anatómica exata da interação coração-pulmão, mas antes como um perfil adicional de movimento, mais rápido e discreto, combinado com a respiração.

\subsection{Tosse e perturbações rápidas}
A tosse é um evento curto, abrupto e menos previsível, provocando, durante uma broncoscopia, a deslocação súbita das vias aéreas, com tendência para a ocorrência de réplicas adjacentes à deslocação principal, alteração temporária do campo visual, variação do fluxo de ar e perda momentânea de orientação. Para quem treina, este tipo de evento é importante porque obriga a recuperar a posição visual e a manter controlo do broncoscópio após uma perturbação.

A tosse é, em si mesma, um reflexo vital complexo, composto por uma fase inspiratória, uma fase compressiva com encerramento da glote e uma fase expiratória explosiva, coordenada por recetores das vias aéreas e por vias nervosas aferentes e eferentes bem descritas na literatura fisiológica \parencite{Andrani2018}. A modelação computacional recente deste fenómeno tem-se debruçado sobre a dinâmica de gotículas e da película mucosa nas vias aéreas superiores durante um evento de tosse, com o objetivo de relacionar as suas características com diferentes patologias pulmonares \parencite{Ilegbusi2025}, o que confirma o interesse científico atual em compreender este movimento para além da sua simples deteção visual durante um procedimento.

\begin{figure}[htbp]
  \centering
  \includegraphics[width=0.72\textwidth]{artigos/tosse_pressao_andrani.png}
  \caption{Variação de pressão e fluxo durante o reflexo da tosse \parencite{Andrani2018}.}
  \label{fig:tosse-pressao-andrani}
\end{figure}

No simulador, a tosse deve por isso ser representada como um evento transitório, com duração reduzida e amplitude superior à do batimento cardíaco, podendo a sua ocorrência ser programada de forma aleatória dentro de limites definidos, o que permite criar cenários de treino mais realistas sem obrigar à reprodução completa da fisiologia da tosse.

\subsection{Interação com o broncoscópio}
A passagem do broncoscópio também influencia a geometria observada, dado que o contacto com as paredes pode causar deformações locais, resistência ao avanço, pequenas oscilações e alterações da orientação da imagem. Esta componente é especialmente importante para o realismo tátil, porque aproxima o simulador da sensação física de navegação no interior das vias aéreas.

Apesar da sua relevância, esta interação mostra-se a mais difícil de implementar, exigindo materiais flexíveis adequados, geometrias calibradas com espessura e resistência adequadas e, idealmente, sensores ou mecanismos capazes de responder ao contacto. Esta dificuldade é também reconhecida na investigação em robótica endoscópica (área mais ampla, da qual a robótica aplicada à broncoscopia faz parte), onde a deteção de força na interação entre instrumento e tecido continua a ser identificada como um dos principais

\begin{figure}[H]
  \centering
  \includegraphics[width=0.70\textwidth]{artigos/forcas_endoscopia_hosseinabadi.png}
  \caption{Forças de contacto em endoscopia robótica \parencite{Hosseinabadi2021}.}
  \label{fig:forcas-endoscopia-hosseinabadi}
\end{figure}

desafios em cirurgia minimamente invasiva \parencite{Hosseinabadi2021}, tendo já sido exploradas soluções como broncoscópios robóticos de balão, especificamente concebidos para procedimentos nas vias aéreas e capazes de aplicar forças controladas sobre as suas paredes sem lesionar o tecido, ainda que este tipo de solução acrescente complexidade mecânica significativa ao sistema \parencite{Li2022}. Esta complexidade torna particularmente difícil a implementação deste tipo de fenómeno num simulador.

\begin{figure}[H]
  \centering
  \includegraphics[width=0.70\textwidth]{artigos/broncoscopio_balao_li.png}
  \caption{Broncoscópio robótico com balão \parencite{Li2022}.}
  \label{fig:broncoscopio-balao-li}
\end{figure}

\subsection{Patologias e deformações locais}
Estenoses, obstruções, variações anatómicas e alterações patológicas modificam a navegação e o aspeto interno das vias aéreas, sendo que a impressão 3D permite criar geometrias específicas para representar alguns destes cenários de forma estática, o que já pode ser útil para treino de reconhecimento anatómico e planeamento de procedimentos. Um exemplo clínico bem documentado deste tipo de fenómeno é a traqueobroncomalácia, uma condição em que o amolecimento da cartilagem traqueal e brônquica provoca colapso dinâmico excessivo das vias aéreas, observado por videobroncoscopia em cerca de 69\% dos asmáticos com doença grave, contra apenas 1,6\% em indivíduos saudáveis \parencite{DalNegro2013}. Também a broncoconstrição ativa tem sido alvo de modelação matemática multiescala, procurando explicar como a complacência não linear da parede brônquica modula a resposta dinâmica das vias aéreas a estímulos broncoconstritores \parencite{Hiorns2014}.

No entanto, a simulação ativa de estenose variável, colapso local ou deformação dinâmica do lúmen exige um conjunto complexo de atuadores distribuídos, materiais compatíveis e uma estratégia de controlo mais complexa.

\begin{figure}[H]
  \centering
  \includegraphics[width=0.66\textwidth]{artigos/modelo_dinamico_fu.png}
  \caption{Modelo impresso com expansão passiva por insuflação \parencite{Fu2023}.}
  \label{fig:modelo-dinamico-fu}
\end{figure}

\section{Síntese crítica}
Esta análise comprova que existem necessidades que nem sempre são satisfeitas em simultâneo. 
Por um lado, os simuladores devem ser suficientemente acessíveis para poderem ser usados em contexto académico e profissional, permitindo treinos repetidos.
Por outro, devem evitar uma representação demasiado estática das vias aéreas, porque a broncoscopia real ocorre num ambiente anatómico em constante movimento.

\begin{figure}[H]
  \centering
  \begin{tikzpicture}[
    box/.style={draw, rounded corners=4pt, align=center, minimum width=3.6cm, minimum height=1.2cm, font=\small, thick},
    arr/.style={-Latex, thick}
  ]
    \node[box, fill=gray!10] (static) {Modelos físicos\\estáticos};
    \node[box, fill=blue!8, right=0.8cm of static] (virtual) {Simulação\\virtual};
    \node[box, fill=orange!12, right=0.8cm of virtual] (dynamic) {Modelos com\\dinamismo parcial};
    \node[box, fill=green!12, below=1.1cm of virtual, minimum width=4.3cm] (gap) {Lacuna identificada:\\árvore brônquica física\\com atuação mecânica};
    \draw[arr] (static.south) -- (gap.west);
    \draw[arr] (virtual.south) -- (gap.north);
    \draw[arr] (dynamic.south) -- (gap.east);
  \end{tikzpicture}
  \caption{Posicionamento da lacuna identificada na literatura revista.}
  \label{fig:sintese-lacuna-literatura}
\end{figure}

Os modelos impressos em 3D respondem bem à primeira necessidade, pois permitem criar geometrias anatómicas de baixo custo e iterar rapidamente o projeto, como demonstram diversos exemplos comerciais e académicos de árvores brônquicas estáticas, alguns já disponíveis há várias décadas \parencite{Pedersen2017,Ho2019,Leba2023,Plan3DChile,TruCorpAirSim,PatUS5597310A}.
Contudo, quando usados isoladamente, tendem a limitar-se à fidelidade geométrica.

Algumas publicações têm procurado atenuar exatamente esta limitação, conferindo algum grau de dinamismo a modelos impressos em 3D e a outras réplicas físicas de vias aéreas, embora nenhuma delas recorra a atuação mecânica da própria estrutura anatómica. Um primeiro exemplo consiste na insuflação controlada de ar ou soro através de uma porta lateral, fazendo com que o próprio modelo impresso se expanda e contraia de forma passiva, uma solução já aplicada a um protótipo de via aérea impresso em 3D capaz de emular movimento fisiológico e patológico \parencite{Fu2023}. Outro trabalho, mais próximo da atuação robótica, recorre a seringas acionadas por uma guia linear motorizada para empurrar ar através de condutas rígidas que representam a árvore brônquica e a traqueia, replicando fisicamente tanto a respiração como a tosse. Contudo, mesmo nesta solução, os brônquios e a traqueia permanecem estruturas rígidas e estáticas, sendo apenas o fluxo de ar no seu interior que é mecanicamente controlado, sem qualquer tentativa de reproduzir o batimento cardíaco \parencite{Giannaccini2017}.

Fora do domínio da broncoscopia, outros trabalhos recorrem a atuação mecânica real para simular a mecânica ventilatória, mas fazem-no ao nível do diafragma ou da parede torácica, não da própria árvore brônquica. É o caso de um simulador respiratório que combina músculos artificiais pneumáticos com um diafragma sintético para insuflar pulmões reais ou artificiais \parencite{Horvath2020}, de um simulador eletromecânico de pulmão que combina componentes poliméricos e tecido pulmonar orgânico para reproduzir uma respiração realista \parencite{Pare2019} e de modelos deformáveis usados em radioterapia para replicar mecanicamente a expansão e a contração de um pulmão, sem qualquer representação da árvore brônquica no seu interior \parencite{Chang2010}. À mesma categoria pertencem ainda simuladores pulmonares controlados por motor de passo ou por pistão servo-controlado, desenvolvidos sobretudo para testar e calibrar ventiladores mecânicos \parencite{Zhao2025,PatUS5975748A}, bem como patentes que atuam mecanicamente a parede torácica de manequins de treino, por bolsas de ar ou por sistemas biela-manivela, para simular respiração ou manobras de reanimação \parencite{PatCN101447141B,PatEP2430628B1}. Outros dispositivos comerciais patenteados representam a árvore brônquica de forma mais completa, incluindo canais para brônquio esquerdo, direito e traqueia, mas destinam-se a treinar procedimentos de drenagem postural com sensores de posição, não a reproduzir movimento fisiológico \parencite{PatCN107346634B}. Uma última categoria de dispositivos recorre efetivamente a atuadores mecânicos reais sobre estruturas anatómicas de treino, mas aplicados à via aérea superior, como a mandíbula, a língua ou a articulação temporomandibular em simuladores de gestão de via aérea e de entubação orotraqueal, nunca chegando a atuar a árvore brônquica em profundidade e as suas bifurcações \parencite{TakanishiWKA}.

Do lado puramente computacional, têm também sido propostos métodos capazes de simular visualmente a deformação respiratória das vias aéreas durante a navegação broncoscópica, por renderização de imagem em tempo real a partir de dados de tomografia computorizada, as imagens médicas que descrevem em cortes transversais a anatomia interna do paciente, mas sem qualquer correspondência física. O resultado é sempre uma simulação gráfica, não um modelo tangível que o formando possa manipular com um broncoscópio real \parencite{DunnBeltran2026}. Já no campo da robótica cirúrgica aplicada à broncoscopia, os poucos manipuladores desenvolvidos destinam-se a operar sobre vias aéreas reais de doentes ou de cadáveres, como ferramentas de tratamento e não como modelos de treino, permanecendo a anatomia sempre estática durante a sua atuação \parencite{Gafford2019}. Mesmo alargando esta análise a protótipos ibero-americanos, como um simulador de broncoscopia com realimentação háptica desenvolvido no México, cuja árvore brônquica nunca chega a ser impressa fisicamente e permanece apenas como um modelo gráfico renderizado \parencite{IPN2012}, a conclusão mantém-se a mesma.

Em síntese, não se encontrou em toda a literatura revista nenhum modelo de árvore brônquica impresso em 3D em que a própria estrutura anatómica seja mecanicamente atuada, por manipuladores robóticos ou mecanismo equivalente, de forma a reproduzir de forma simultânea os principais fenómenos fisiológicos observáveis durante a broncoscopia. É precisamente esta lacuna que o presente projeto procura preencher.

\begin{figure}[H]
  \centering
  \begin{tikzpicture}[
    block/.style={draw, rounded corners=4pt, align=center, minimum height=1.0cm, font=\small, thick},
    arr/.style={-Latex, thick}
  ]
    \node[block, fill=blue!8, minimum width=3.1cm] (anat) {Fidelidade\\anatómica};
    \node[block, fill=orange!12, right=0.55cm of anat, minimum width=3.1cm] (mov) {Movimento\\fisiológico};
    \node[block, fill=green!12, right=0.55cm of mov, minimum width=3.1cm] (cost) {Acessibilidade\\e repetição};
    \node[block, fill=gray!10, below=1.0cm of mov, minimum width=5.6cm] (proj) {Objetivo do projeto\\modelo físico dinâmico};
    \draw[arr] (anat.south) -- (proj.north west);
    \draw[arr] (mov.south) -- (proj.north);
    \draw[arr] (cost.south) -- (proj.north east);
  \end{tikzpicture}
  \caption{Critérios reunidos na proposta desenvolvida.}
  \label{fig:sintese-criterios-projeto}
\end{figure}

A componente dinâmica, mesmo que simplificada, acrescenta valor ao treino, pois obriga o utilizador a lidar com deslocações por vezes imprevisíveis e com a perda parcial do enquadramento, devido a perturbações temporárias.

Assim, o objetivo deste projeto encontra-se em combinar uma árvore brônquica física, obtida por impressão 3D, com uma arquitetura de atuação capaz de reproduzir de forma controlada os fenómenos anatómicos dinâmicos identificados ao longo deste capítulo, respondendo a uma lacuna que a literatura revista ainda não preenche.
